\centerline{\bf \Huge Двойные отношения}

\begin{flushright}
        \scriptsize{Всем доброе утро!\\ Сегодня я что-то заблудился на жизненном пути...\\ Какаши Хатаке}
\end{flushright}

\textbf{Определение.} Двойным отношением четвёрки различных(!) точек $A, B, C, D$, лежащих на одной прямой, называется число

$$
(A, B ; C, D)=\frac{\overrightarrow{A C}}{\overrightarrow{B C}}: \frac{\overrightarrow{A D}}{\overrightarrow{B D}}
$$


Если формулировать словами, то двойное отношение четвёрки $A, B, C, D$ - это отношение отношений, в которых точки $C$ и $D$ делят отрезок $A B$.

\problem{1}
Докажите, что по трём точкам $A, B, C$ и двойному отношению $(A, B, C, D)$ однозначно восстанавливается точка $D$.

\problem{2}
Чему равно двойное отношение $(A, B, C, D)$, если $D$ - бесконечно удалённая?

\textbf{Определение.} Двойным отношением упорядоченной четверки прямых $a, b, c, d$, называется величина

$$
(a, b, c, d)=\frac{\sin \angle(\bar{a}, \bar{c})}{\sin \angle(\bar{b}, \bar{c})}: \frac{\sin \angle(\bar{a}, \bar{d})}{\sin \angle(\bar{b}, \bar{d})},
$$


где $\bar{a}, \bar{b}, \bar{c}, \bar{d}$ --- произвольные векторы, направленные вдоль прямых $a, b, c, d$.

\problem{3}
(а) Пусть прямые $a, b, c, d$ пересекаются в одной точке $O$, а прямая $l$ пересекает эти прямые в точках $A, B, C, D$ соответственно. Докажите, что $(A, B, C, D)=(a, b, c, d)$. Выведите отсюда аналог задачи 1 для прямых, а так же что двойное отношение четвёрки точек сохраняется при центральной проекции.

\noindent(б) Прямые $\ell_1$ и $\ell_2$ пересекаются в точке $A$. На прямой $\ell_1$ отмечены точки $B_1, C_1, D_1$, а на прямой $\ell_2$ - точки $B_2, C_2, D_2$. Докажите, что прямые $B_1 B_2, C_1 C_2, D_1 D_2$ конкурентны (то есть пересекаются в одной точке или параллельны) тогда и только тогда, когда $\left(A, B_1, C_1, D_1\right)=\left(A, B_2, C_2, D_2\right)$.

\textbf{Определение.} Четверка точек (прямых) называется гармонической, если их двойное отношение равно -1.

\newpage

\hspace{3mm}

\problem{4}
Известно, что $(A, B ; C, D)=\lambda$. Найдите $(B, A ; C, D),(A, B ; D, C)$ и $(A, C ; B, D)$. Осознайте, что с помощью этого можно найти двойное отношение для любой перестановки точек.

\problem{5} Пусть $A, B, C, D$ - различные точки, и $(A, B, C, D)=(B, A, C, D)$. Докажите, что $(A, B, C, D)=-1$.

\problem{6} Докажите, что следующие четверки гармонические.

(a) $\left(A, B, M, P_{\infty}\right)$, где $M$ - середина отрезка $A B$, а $P_{\infty}$ - бесконечно удаленная точка вдоль прямой $A B$.

(b) $(B, C, M, N)$, где $M$ и $N$ - основания внутренней и внешней биссектрис треугольника $A B C$ с основанием $B C$.

(c) $(A, B, X, Y)$, где $X$ и $Y$ центры непересекающихся окружностей разного радиуса, $A$ и $B$ - точки пересечения общих внешних и внутренних касательных.

\problem{7} 
\textbf{Теорема о полном четырёхстороннике.} Продолжения сторон $A B$ и $C D$ четырехугольника $A B C D$ пересекаются в точке $E$, продолжения $B C$ и $A D$ - в точке $F$, прямые $A C$ и $B D$ пересекают $E F$ в точках $M$ и $N$. Докажите, что $(E, F, M, N)=-1$.

\problem{8}
\textbf{Теорема Паппа.} На двух прямых располагаются по три точки $A, B, C$ и $A^{\prime}, B^{\prime}$ $C^{\prime}$. Тогда точки пересечения пар прямых $A B^{\prime}$ и $A^{\prime} B, B C^{\prime}$ и $B^{\prime} C, C A^{\prime}$ и $C^{\prime} A$ лежат на одной прямой.

\problem{9}
\textbf{Лемма об изогоналях.} Внутри угла $A O B$ проведены лучи $O D$ и $O C$, симметричные относительно биссектрисы этого угла. Если $M$ - точка пересечения $A D$ и $B C$, а $N$ - точка пересечения $B D$ и $A C$, то лучи $O N$ и $O M$ также симметричны относительно биссектрисы угла $A O B$.